\clearpage\thispagestyle{empty}
\pdfbookmark[0]{Quellen zum Bericht über Lindenaus Untersuchungen zum Sonnendurchmesser}{lindenausolar1809-sources}
\begin{center}\textbf{Quellen zum Bericht über Lindenaus Untersuchungen zum Sonnendurchmesser}\end{center}
\noindent Diese KI-übersetzte Fassung in modernem Deutsch folgt Carl Friedrich Gauß, \href{https://gdz.sub.uni-goettingen.de/id/PPN236007467}{\emph{Werke}, Band~VI} (Göttingen, 1874), S.~517–520. Der Bericht erschien zuerst ohne Verfassernamen unter Göttingen in den \emph{Göttingischen gelehrten Anzeigen}, Stück~113, 17. Juli 1809. Der \href{https://images.p.sub.uni-goettingen.de/iiif/image/gdz:PPN319721507_1809:00001125/full/1650,/0/default.jpg}{Anfang} berichtet Ergebnisse aus Bernhard von Lindenaus Brief vom 8. Juni an Gauß. Am Ende folgt ein Hinweis auf Sternbedeckungen durch den Mond aus dem Verzeichnis von Herrn von Wisniewski, das Heinrich Wilhelm Olbers weitergab. Die Kopfzeile der Werkausgabe nennt S.~1121–1126. Der Hinweis endet in der Zeitschrift auf S.~1127.
\par\smallskip\noindent Im Erstdruck steht auf \href{https://images.p.sub.uni-goettingen.de/iiif/image/gdz:PPN319721507_1809:00001129/full/1650,/0/default.jpg}{S.~1125} $0{,}9631$ und $117/116$, in der Werkausgabe $0{,}09631$ und $17/116$. Die Übertragung folgt der Werkausgabe. In seinem Bericht vom 16. Juni 1810 (\emph{Werke}, Band~VI, S.~525) bezeichnet Gauß den ersten dieser Werte im Erstdruck als Druckfehler. Auf \href{https://images.p.sub.uni-goettingen.de/iiif/image/gdz:PPN319721507_1809:00001131/full/1650,/0/default.jpg}{S.~1127} steht in der Zeitschrift die Tabelle der Sternbedeckungen vor dem Schlusssatz. In der Werkausgabe stehen an ihrer Stelle zwei Striche. Übersetzung und Satz mit KI erstellt.
\clearpage
